\chapter{Bøyingslære}

Dette kapittelet gir en innføring i bøying av skoltesamiske ord med tilhørende øvingsoppgaver. Fokuset er på selve bøyinga av orda. For en beskrivelse av hvordan de ulike formene skal brukes, se kapitla om kasusbruk og setningslære.

I første delkapittel beskrives oppbygninga av de skoltesamiske orda og hvordan de kan deles i uike grupper. Andre delkapittel gir en grundig innføring i stadieveksling og hvordan denne brukes i bøying av enstava substantiv, adjektiv og tallord. Tredje delkapittel beskriver vokalveksling, også dette med eksempel og oppgaver knytta til samme slags ord. Fjerde delkapittel beskriver verbbøyinga. Femte delkapittel beskriver bøyinga av alle nomen, dvs. substantiv, adjektiv, tallord og pronomen [INNDELING: allment gruppe for gruppe, spesielt for adjektiv, henvisning til syntakskapittel for fraser med tallord og/eller pronomen + substantiv]. I sjette kapittel behandles eiendomsendelsene.

\section{Oppbygging og inndeling av skoltesamiske ord} \label{inndeling}



\section{Stadieveksling}

[Grad og stadium]

\subsection{Typen YXX : YX}

I denne typen ord består konsonantsentrumet av to ulike konsonanter, her representert ved stor Y og stor X. Dette er ord som Mäʹrjj 'Maria', põlvv 'sky', äldd 'reinsimle' og mueʹrjj 'bær'. Vokalene \emph{i} og \emph{u} regnes i stadievekslingsøyemed som konsonanter når de kommer etter en annen vokal, slik som i orda leiʹbb 'brød' og sauʒʒ 'sau', og hører også til denne gruppa.

Svakt stadium av denne typen dannes ved å forlenge vokalen og forkorte den andre konsonanten på dette viset:

Mäʹrjj : Määʹrj

põlvv : põõlv

äldd : ääld

mueʹrjj : mueʹrj

leiʹbb : leeiʹb

sauʒʒ : saauʒ

Merk at forskjellen på lang og kort diftong ikke markeres i skrift.

Dette er nompl accsg gensg, gi eksempel.

Oppgaver

\subsubsection{Undergruppa Ygg : Yǥ}

I de orda der den andre konsonanten er \emph{gg}, endres denne til \emph{ǥ} (g med strek) i svakt stadium i tillegg til de vanlige lengdeendringene, f.eks. argg 'hverdag' : aarǥ, ...

Dette gjelder ikke om den første konsonanten er \emph{ŋ}, f.eks. siâŋgg 'seng' : siâŋg.

Oppgaver

\subsubsection{Undergruppa Yǧǧ : Yj}

I de orda der den andre konsonanten er \emph{ǧǧ}, endres denne til \emph{j} i svakt stadium i tillegg til de vanlige lengdeendringene, f.eks.  âʹlǧǧ 'gutt' : ââʹlj, jueʹlǧǧ 'fot, bein' : jueʹlj, ...

Dette gjelder ikke om den første konsonanten er \emph{ŋ}, f.eks. šäʹŋǧǧ 'pirog' : šääʹŋǧ.

Oppgaver

\subsubsection{Undergruppa hXX : uX}

I de orda der den første konsonanten er \emph{h}, endres denne til \emph{u} i svakt stadium i tillegg til de vanlige lengdeendringene, f.eks.  

Dette gjelder ikke om den andre konsonanten er \emph{k}, f.eks. vihkk 'notisbok' : viihk.

Merk at det likevel aldri vil bli tre u-er etter hverandre selv om den første vokalen er u. I disse tilfellene blir vokalen forlenga og h faller bort, f.eks. uhss 'dør' : uus.

Oppgaver

\subsubsection{Undergruppa jXX : X og vXX : X}

Etter kort \emph{i} brukes ikke \emph{i} som en del av konsonantsentrumet, siden dette vil se ut som en lang \emph{ii}. I steden skrives \emph{j}, f.eks. sijdd '(skolte)bygd'. Når vokalen forlenges i svakt stadium, faller \emph{j} bort, altså sijdd : siid.

Etter kort \emph{u} brukes ikke \emph{u} som en del av konsonantsentrumet, siden dette vil se ut som en lang \emph{uu}. I steden skrives \emph{v}, f.eks. vuʹvdd 'område'. Når vokalen forlenges i svakt stadium, faller \emph{v} bort, altså vuʹvdd : vuuʹd.

Oppgaver

Oppsummerende oppgaver

\subsection{Typen XˈX : XX}

I denne typen består konsonantsentrumet av to like konsonanter, enten bb, dd, gg, ʒʒ, ǯǯ eller llj. Dette er ord som låʹdd 'fugl', kõʒʒ 'negl' og villj 'bror'. I svakt stadium, grad II, blir vokalen forlenga mens konsonantsentrumet forblir uendra, f.eks. låʹdd : lååʹdd, kõʒʒ : kõõʒʒ, villj: viillj.

Merk at siden diftonglengde ikke merkes i skrift, vil ekstra sterk og sterk grad se like ut selv om de uttales ulikt. Med ordboksortografi ser man forskjellen f.eks. vueʹǯˈǯ 'kjøtt' : vueʹǯǯ.

Oppgaver

\subsection{Typen XˈX : XX : X}

I denne typen består konsonantsentrumet av to like, stemte konsonanter. Disse er đđ, jj, ll, mm, nn, nnj, ŋŋ og vv.

Denne gruppa har tre ulike grader: III, II og I. I grad III er det kort vokal og lang konsonant, i grad II lang vokal og lang konsonant og i grad I lang vokal og kort konsonant.

Disse veksler vanligvis slik at dersom grunnforma med sterkt stadium har grad III, har svakt stadium grad II, og dersom grunnforma med sterkt stadium har grad II, har svakt stadium grad I. Eksempler

Grad III : grad II:

kåʹll 'gull' : kååʹll

cuʹmm 'kyss' : cuuʹmm

käll 'panne (kroppsdel)' : kääll

Grad II : grad I

njuuʹnn 'nese' : njuuʹn

luumm 'plomme (frukt)' : luum

mââʹnnj 'svigerinne' : mââʹnj

ǩiõll 'språk' : ǩiõl

Mange ord med â-stamme (se tidligere underkapittel \ref{inndeling}) "hopper over" en grad, dvs. at i sterkt stadium har de grad III, med kort vokal og lang konsonant, og i svakt stadium har de grad I, med lang vokal og kort konsonant. Eksempler:

võrr 'blod' : võõr

toll 'bål' : tool

Oppgaver (merk ord som hopper over et stadium)

\subsection{Typen XˈX : XX : (Z)Z}

I denne typen består konsonantsentrumet av to like, ustemte konsonanter. Disse er cc, čč, kk/ǩǩ, pp, tt, ss og šš.

Denne typen minner mye om typen XˈX : XX : X. Forskjellen er at grad I kjennetegnes av at konsonantene blir stemte, ikke nødvendigvis kortere. Konsonantene tt og pp blir både stemte og korte.

Stemthetsvekslingene til grad I er som følger:

cc : ʒʒ

čč : ǯǯ

kk : ǥǥ

ǩǩ : jj

pp : v

tt : đ

ss : zz

šš : žž

Eksempler på vekslingene:

Grad III > grad II:

kopp 'kopp' : koopp

vacc 'vott' : vaacc

äʹšš 'sak' : ääʹšš

Grad II > grad I:

čääʹcc 'vann' : čääʹʒʒ

laakk 'tak' : laaǥǥ

šââʹǩǩ 'gris' : šââʹjj

sieʹss 'tante (fars søster)' : sieʹzz

Også i denne gruppa kan mange ord med â-stamme (se tidligere underkapittel \ref{inndeling}) "hoppe over" en grad, dvs. at i sterkt stadium har de grad III, med kort vokal og lang, ustemt konsonant, og i svakt stadium har de grad I, med lang vokal og stemt konsonant. Eksempler:

jokk 'elv' : jooǥǥ

kuss 'ku: kuuzz

Oppgaver

\subsection{Vokalveksling knytt til stadievekslinga}

Diftongene uẹʹ\~uäʹ og iẹʹ\~eäʹ veksler etter hva slags konsonantsentrum som følger. Om konsonantsentrumet er av typen XˈX, X, Z eller YXX, blir diftongen uäʹ / eäʹ. Om konsonantsentrumet derimot er av typen XX, ZZ eller YX, blir diftongen uẹʹ / iẹʹ

Eksempel:

jeäʹnˈn 'mor' : jiẹʹnn

heäʹrvv 'pynt' : hiẹʹrv

nueʹđđ 'bør' : nuäʹđ

Merk at dette handler om den faktiske lengden på den påfølgende konsonanten, ikke graden, som man ser i det følgende eksemplet:

ǩiẹʹss 'sommer' : ǩiẹʹzz

E med prikk under brukes i disse uẹʹ- og iẹʹ-diftongene for å skille dem fra de ueʹ- og ieʹ-diftongene som ikke endres avhengig av hvilket konsonantsentrum som følger. Slike som ikke veksler, finnes f.eks. i orda mieʹlǩǩ 'melk' : mieʹlǩ og kueʹll 'fisk' : kueʹl.


\subsection{Omvendt stadieveksling}

homonymi cc:ʒʒ ʒʒ:ʒʒ

\section{Vokalveksling}

\section{Verbbøying}

\section{Nomenbøying}

\section{Eiendomsendelser}